% GENERATED FILE. Edit canonical lessons, then run npm run generate:books.
% canonical-source-hash: fnv1a64:490128eedfd536ad
% canonical-lessons: MW-C217-laghno, MW-C217-charhno, MW-C217-bojh, MW-C217-ladno, MW-C217-thairno

\chapter{Over the River}
\label{ch:mw-over-the-river}
\hlchaptermodality{217}

\begin{chapteropening}
\textbf{By the end of this chapter:} \emph{I can say I will cross the river, load the camel and climb on.}

Complete the last lesson of chapter 217: I can say I will cross the river, load the camel and climb on.
\end{chapteropening}

\section[\texorpdfstring{\mw{लांघणो} (lā̃ghṇo)}{लांघणो (lā̃ghṇo)}]{\mw{लांघणो} (lā̃ghṇo) — to cross over; to leap over}

\label{lesson:MW-C217-laghno}



\begin{quote}
Before the new one: say the Marwadi for a traveller, then the Marwadi for a camp.
\end{quote}

\subsection*{You'll want to know: \mw{लांघणो}}
\textbf{\mw{लांघणो}} — \emph{lā̃ghṇo} — \textquotedblleft{}to cross over; to leap over\textquotedblright{}.

\begin{quote}
\mw{हूं} \mw{काल} \mw{नदी} \mw{लांघसूं।}
\end{quote}

\begin{quote}
\emph{hū̃ kāl nadī lā̃ghsū̃.}
\end{quote}

\begin{quote}
Tomorrow I'll cross the river.
\end{quote}

Past (I …): \mw{म्हैं} \mw{नदी} \mw{लांघी}, \emph{mhaĩ nadī lā̃ghī}.

\begin{cousinweb}
\textbf{\mw{लांघणो}} (\emph{lā̃ghṇo}) is inherited from Sanskrit \emph{laṅghayati}, \textquotedblleft{}leaps over\textquotedblright{}; CDIAL lists Hindi \emph{lā̃ghnā}, \textquotedblleft{}to jump over\textquotedblright{}, under it. CDIAL gives no Marwari form in this entry (CDIAL 10905).
\end{cousinweb}

\begin{grammarlens}[title={What you've built}]
One more word for saying what comes next.
\end{grammarlens}

\subsection*{Your turn}
\begin{itemize}
\raggedright
  \item \emph{Say it:} \emph{lā̃ghṇo}
  \item \emph{Say it:} \emph{lā̃ghṇo}, once more
  \item \emph{Say it:} \emph{hū̃ kāl nadī lā̃ghsū̃.} — Tomorrow I'll cross the river.
  \item \emph{Recall:} say the Marwadi for a traveller, then the Marwadi for a camp, then say \emph{lā̃ghṇo} again
\end{itemize}

\begin{tcolorbox}[breakable,colback=teal!4,colframe=teal!35!black,title={Before you move on}]
What does \mw{लांघणो} mean? (\textquotedblleft{}To cross over; to leap over\textquotedblright{}.) Say it once more.
\end{tcolorbox}
\section[\texorpdfstring{\mw{चढ़णो} (chaṛhṇo)}{चढ़णो (chaṛhṇo)}]{\mw{चढ़णो} (chaṛhṇo) — to climb; to get on (an animal or a vehicle)}

\label{lesson:MW-C217-charhno}



\begin{quote}
Before the new one: say the Marwadi for a camp, then the Marwadi for to cross over.
\end{quote}

\subsection*{You'll want to know: \mw{चढ़णो}}
\textbf{\mw{चढ़णो}} — \emph{chaṛhṇo} — \textquotedblleft{}to climb; to get on (an animal or a vehicle)\textquotedblright{}.

\begin{quote}
\mw{हूं} \mw{काल} \mw{ऊंट} \mw{माथै} \mw{चढ़सूं।}
\end{quote}

\begin{quote}
\emph{hū̃ kāl ū̃ṭ māthai chaṛhsū̃.}
\end{quote}

\begin{quote}
Tomorrow I'll climb up onto the camel.
\end{quote}

Past (he …): \mw{वो} \mw{ऊंट} \mw{माथै} \mw{चढ़्यो}, \emph{vo ū̃ṭ māthai chaṛhyo}.

\begin{cousinweb}
\textbf{\mw{चढ़णो}} (\emph{chaṛhṇo}) is inherited from a reconstructed \emph{*caḍhati}, \textquotedblleft{}rises\textquotedblright{}; CDIAL lists Marwari \emph{caṛhṇo}, \textquotedblleft{}to climb\textquotedblright{} (CDIAL 4578).
\end{cousinweb}

\begin{grammarlens}[title={What you've built}]
One more word for saying what comes next.
\end{grammarlens}

\subsection*{Your turn}
\begin{itemize}
\raggedright
  \item \emph{Say it:} \emph{chaṛhṇo}
  \item \emph{Say it:} \emph{chaṛhṇo}, once more
  \item \emph{Say it:} \emph{hū̃ kāl ū̃ṭ māthai chaṛhsū̃.} — Tomorrow I'll climb up onto the camel.
  \item \emph{Recall:} say the Marwadi for a camp, then the Marwadi for to cross over, then say \emph{chaṛhṇo} again
\end{itemize}

\begin{tcolorbox}[breakable,colback=teal!4,colframe=teal!35!black,title={Before you move on}]
What does \mw{चढ़णो} mean? (\textquotedblleft{}To climb; to get on (an animal or a vehicle)\textquotedblright{}.) Say it once more.
\end{tcolorbox}
\section[\texorpdfstring{\mw{बोझ} (bojh)}{बोझ (bojh)}]{\mw{बोझ} (bojh) — a burden; a heavy load}

\label{lesson:MW-C217-bojh}



\begin{quote}
Before the new one: say the Marwadi for to cross over, then the Marwadi for to climb.
\end{quote}

\subsection*{You'll want to know: \mw{बोझ}}
\textbf{\mw{बोझ}} — \emph{bojh} — \textquotedblleft{}a burden; a heavy load\textquotedblright{}.

\begin{quote}
\mw{ऊंट} \mw{माथै} \mw{घणो} \mw{बोझ} \mw{है।}
\end{quote}

\begin{quote}
\emph{ū̃ṭ māthai ghaṇo bojh hai.}
\end{quote}

\begin{quote}
The camel is carrying a heavy load.
\end{quote}

A masculine noun.

\begin{cousinweb}
\textbf{\mw{बोझ}} (\emph{bojh}) is inherited from Sanskrit \emph{vahya}, \textquotedblleft{}to be carried\textquotedblright{}; CDIAL lists Hindi \emph{bojh}, \textquotedblleft{}a load\textquotedblright{}, under it. CDIAL gives no Marwari form in this entry (CDIAL 11465).
\end{cousinweb}

\begin{grammarlens}[title={What you've built}]
One more word for saying what comes next.
\end{grammarlens}

\subsection*{Your turn}
\begin{itemize}
\raggedright
  \item \emph{Say it:} \emph{bojh}
  \item \emph{Say it:} \emph{bojh}, once more
  \item \emph{Say it:} \emph{ū̃ṭ māthai ghaṇo bojh hai.} — The camel is carrying a heavy load.
  \item \emph{Recall:} say the Marwadi for to cross over, then the Marwadi for to climb, then say \emph{bojh} again
\end{itemize}

\begin{tcolorbox}[breakable,colback=teal!4,colframe=teal!35!black,title={Before you move on}]
What does \mw{बोझ} mean? (\textquotedblleft{}A burden; a heavy load\textquotedblright{}.) Say it once more.
\end{tcolorbox}
\section[\texorpdfstring{\mw{लादणो} (lādṇo)}{लादणो (lādṇo)}]{\mw{लादणो} (lādṇo) — to load (an animal or a cart)}

\label{lesson:MW-C217-ladno}



\begin{quote}
Before the new one: say the Marwadi for to climb, then the Marwadi for a burden.
\end{quote}

\subsection*{You'll want to know: \mw{लादणो}}
\textbf{\mw{लादणो}} — \emph{lādṇo} — \textquotedblleft{}to load (an animal or a cart)\textquotedblright{}.

\begin{quote}
\mw{हूं} \mw{काल} \mw{ऊंट} \mw{माथै} \mw{बोझ} \mw{लादसूं।}
\end{quote}

\begin{quote}
\emph{hū̃ kāl ū̃ṭ māthai bojh lādsū̃.}
\end{quote}

\begin{quote}
Tomorrow I'll put the load on the camel.
\end{quote}

Past (I …): \mw{म्हैं} \mw{बोरो} \mw{लाद्यो}, \emph{mhaĩ boro lādyo}.

\begin{cousinweb}
\textbf{\mw{लादणो}} (\emph{lādṇo}) is inherited from Sanskrit \emph{lardayati}, \textquotedblleft{}loads\textquotedblright{}; CDIAL lists the Old Marwari verb \emph{lādaï}, \textquotedblleft{}loads\textquotedblright{} (CDIAL 10966).
\end{cousinweb}

\begin{grammarlens}[title={What you've built}]
One more word for saying what comes next.
\end{grammarlens}

\subsection*{Your turn}
\begin{itemize}
\raggedright
  \item \emph{Say it:} \emph{lādṇo}
  \item \emph{Say it:} \emph{lādṇo}, once more
  \item \emph{Say it:} \emph{hū̃ kāl ū̃ṭ māthai bojh lādsū̃.} — Tomorrow I'll put the load on the camel.
  \item \emph{Recall:} say the Marwadi for to climb, then the Marwadi for a burden, then say \emph{lādṇo} again
\end{itemize}

\begin{tcolorbox}[breakable,colback=teal!4,colframe=teal!35!black,title={Before you move on}]
What does \mw{लादणो} mean? (\textquotedblleft{}To load (an animal or a cart)\textquotedblright{}.) Say it once more.
\end{tcolorbox}
\section[\texorpdfstring{\mw{ठैरणो} (ṭhairṇo)}{ठैरणो (ṭhairṇo)}]{\mw{ठैरणो} (ṭhairṇo) — to halt; to stay for a while}

\label{lesson:MW-C217-thairno}



\begin{quote}
Before the new one: say the Marwadi for a burden, then the Marwadi for to load.
\end{quote}

\subsection*{You'll want to know: \mw{ठैरणो}}
\textbf{\mw{ठैरणो}} — \emph{ṭhairṇo} — \textquotedblleft{}to halt; to stay for a while\textquotedblright{}.

\begin{quote}
\mw{हूं} \mw{परसूं} \mw{तांई} \mw{गांव} \mw{मांय} \mw{ठैरसूं।}
\end{quote}

\begin{quote}
\emph{hū̃ parsū̃ tā̃ī gā̃v mā̃y ṭhairsū̃.}
\end{quote}

\begin{quote}
I'll stop in the village until the day after tomorrow.
\end{quote}

Past (he …): \mw{वो} \mw{गांव} \mw{मांय} \mw{ठैर्यो}, \emph{vo gā̃v mā̃y ṭhairyo}.

\begin{cousinweb}
\textbf{\mw{ठैरणो}} (\emph{ṭhairṇo}) goes back to a reconstructed \emph{*stabhira}, \textquotedblleft{}fixed, firm\textquotedblright{}; CDIAL lists Hindi \emph{ṭhaharnā}, \textquotedblleft{}to stop\textquotedblright{}, and a Gujarati verb \textquotedblleft{}to stay\textquotedblright{} under it. CDIAL gives no Marwari form in this entry (CDIAL 13680).
\end{cousinweb}

\begin{grammarlens}[title={What you've built}]
One more word for saying what comes next.
\end{grammarlens}

\subsection*{Your turn}
\begin{itemize}
\raggedright
  \item \emph{Say it:} \emph{ṭhairṇo}
  \item \emph{Say it:} \emph{ṭhairṇo}, once more
  \item \emph{Say it:} \emph{hū̃ parsū̃ tā̃ī gā̃v mā̃y ṭhairsū̃.} — I'll stop in the village until the day after tomorrow.
  \item \emph{Recall:} say the Marwadi for a burden, then the Marwadi for to load, then say \emph{ṭhairṇo} again
\end{itemize}

\begin{tcolorbox}[breakable,colback=teal!4,colframe=teal!35!black,title={Before you move on}]
What does \mw{ठैरणो} mean? (\textquotedblleft{}To halt; to stay for a while\textquotedblright{}.) Say it once more.
\end{tcolorbox}
